% META: {"id": "1NSI-WEB-CO-017", "chapitre": "1NSI-WEB-IHM", "type_objet": "correction", "exercice_ref": "1NSI-WEB-EX-005", "capacites": ["1NSI-WEB-IHM-C3"], "sources_inspiration": [], "status": "approved", "capacites_codes": ["C3"], "parcours": 1, "fichier_exercice": "chapitres/1NSI-WEB-IHM/exercices/1NSI-WEB-EX-017.tex"}
% BEGIN-VERIFY
% from urllib.parse import urlencode
% params = {"destination": "Djerba", "voyageurs": 2, "tri": "prix_croissant"}
% url = "https://voyages.fr/recherche?" + urlencode(params)
% assert url == "https://voyages.fr/recherche?destination=Djerba&voyageurs=2&tri=prix_croissant"
% END-VERIFY
\begin{corrige}{1NSI-WEB-EX-005}
\textbf{1.}
\begin{python}
from urllib.parse import urlencode

params = {"destination": "Djerba", "voyageurs": 2, "tri": "prix_croissant"}
url = "https://voyages.fr/recherche?" + urlencode(params)
print(url)
# https://voyages.fr/recherche?destination=Djerba&voyageurs=2&tri=prix_croissant
\end{python}

\textbf{2.} Oui : avec GET, les paramètres de recherche font partie intégrante de l'URL.
Mettre cette URL en favori (ou la partager) permet de retrouver exactement la même
recherche plus tard, car tous les critères sont encodés dans l'adresse elle-même.

\textbf{3.} Non, ce n'est pas justifié ici : les données transmises (destination,
nombre de voyageurs, critère de tri) ne sont \textbf{pas confidentielles} — ce sont des
critères de recherche publics, pas des mots de passe ni des informations sensibles. GET est
adapté, et présente même un avantage (URL partageable/favorisable) que POST n'offre pas.
\end{corrige}
